% AC-S50: finite arithmetic multiplicity reconstruction; separate from the ground line.
\section{Multiplicity-sensitive finite arithmetic compression}
\label{sec:spectral-arithmetic-multiplicity}

This is a separate construction from the lowest eigenline of the
localized Weil operator.  It uses the positive, unshifted Weil metric.
The statement below is conditional on \(A\), and does not prove the
ratio (RZ.3), eventual simple-even ground behavior, or a uniform
second-eigenvalue lower bound.  It also supplies no effective rate for
the final convergence.

Use
\[
 \widehat f(t)=\int_{\mathbb R}f(x)e^{-itx}\,dx,\qquad
 \Xi(t)=\xi(1/2+it).
 \tag{MR.1}
\]
Under \(A\), write \(\Gamma\) for the distinct ordinates of the
nontrivial zeros and \(m_\gamma\) for their multiplicities.  The
global Weil completion \(H\) has the sampling isometry
\[
 q(f,g)=\sum_{\gamma\in\Gamma}
       m_\gamma\widehat f(\gamma)\overline{\widehat g(\gamma)},
 \qquad H\simeq\ell^2(\Gamma,m).
 \tag{MR.2}
\]
The equality of pairings in (MR.2) is the ordinary Weil explicit
formula under \(A\).  Step 0 below proves the surjectivity of the
completion map directly; no separate global spectral-realization
theorem is assumed.  Suzuki, \emph{Spectral convergence associated
with Weil's quadratic form}, arXiv:2610.32704v1, Sections 2--3,
uses the same completed realization.  In the present convention
\(D=-i\,d/dx\) is multiplication by \(\gamma\), and
\[
 C=D^{-1},\qquad (C F)(\gamma)=\gamma^{-1}F(\gamma)
 \tag{MR.3}
\]
is compact, selfadjoint and reflection-odd.  A multiplicity
\(m_\gamma>1\) changes the norm at that coordinate; the eigenvalue
\(1/\gamma\) of \(C\) still has a one-dimensional eigenspace.
There is no ordinate zero: by DLMF 25.2.3 at \(s=1/2\),
\[
 (1-\sqrt2)\zeta(1/2)
   =\sum_{r\ge1}\bigl((2r-1)^{-1/2}-(2r)^{-1/2}\bigr)>0.
\]
The grouped alternating series converges and is strictly positive,
so \(\zeta(1/2)\ne0\); the remaining factors in \(\xi(1/2)\)
are finite and nonzero.  The positive distinct ordinates
will be denoted \(0<\gamma_1<\gamma_2<\cdots\).
These identifications are used to prove convergence, not as input
to the finite construction.

\paragraph{The fixed test spaces.}
For \(N\ge1\), let \(V_N\) be the complex vector space of continuous
piecewise-linear functions on \([-N,N]\), with mesh \(2^{-N}\),
which vanish at the endpoints and have integral zero.  Extend
them by zero to \(\mathbb R\).  These spaces are nested and
reflection-stable.  Their dimensions are
\[
 d_N=2N\,2^N-2.
 \tag{MR.4}
\]
Fix a real rational basis separated into its even and odd parts,
using one fixed rational elimination rule.  The corresponding
reflection matrix \(R_N\) is diagonal with entries \(1,-1\).

For a compact mean-zero function put
\[
 (\mathcal P f)(x)=i\int_{-\infty}^x f(y)\,dy.
 \tag{MR.5}
\]
It is compactly supported in the same interval.  For \(f\in V_N\)
it is piecewise quadratic and
\(\widehat{\mathcal P f}(t)=\widehat f(t)/t\), with the value at
zero interpreted by continuity.  Thus it realizes \(Cf\) in
(MR.3).

Fix, independently of \(N\), the real even probe
\[
 \phi(x)=(1-|x|)_++(\sqrt2-|x|)_+.
 \tag{MR.6}
\]
Its transform is
\[
 \widehat\phi(t)=
 \begin{cases}
 [\,2(1-\cos t)+2(1-\cos(\sqrt2\,t))\,]/t^2,&t\ne0,\\
 3,&t=0.
 \end{cases}
 \tag{MR.7}
\]
It is strictly positive for every real \(t\).  Indeed both
summands in its numerator are nonnegative, and their simultaneous
vanishing away from zero would make \(\sqrt2\) rational.
The probe belongs to every local \(H^1_0\) space whose interval
strictly contains \([-\sqrt2,\sqrt2]\), hence to \(H\).

\paragraph{Exact finite matrices.}
Let \(\Pi_N\) denote the Weil-orthogonal projection from \(H\) onto
\(V_N\).  For the chosen basis \(b_1,\ldots,b_{d_N}\), form
\[
 \begin{aligned}
 G_N&=(q(b_j,b_i))_{i,j},&
 H_N&=(q(\mathcal P b_j,b_i))_{i,j},\\
 c_N&=(q(\phi,b_i))_i,&
 S_N&=G_N^{-1/2}H_NG_N^{-1/2},\qquad
 a_N=G_N^{-1/2}c_N .
 \end{aligned}
 \tag{MR.8}
\]
The symbol \(a_N\) here is a probe-coordinate vector, not a
localized Weil eigenvalue.  Below it is always written with
its subscript \(N\); the bottom eigenvalues remain \(a_1(L)\).

The Gram matrix \(G_N\) is positive definite by the local strict
positivity under \(A\), already recorded in AC-S03.  It is real and block diagonal by parity.
The matrix \(H_N\) is Hermitian, purely imaginary and block
off-diagonal.  Consequently \(S_N\) is Hermitian and
\(R_NS_NR_N=-S_N\); \(a_N\) is real and even.
The map sending a Euclidean coordinate vector \(v\) to
\(\sum_j(G_N^{-1/2}v)_jb_j\) is an isometry onto \(V_N\).
Under this isometry \(S_N\) is exactly the Weil-orthogonal
compression \(\Pi_NC\Pi_N\), and \(a_N\) represents \(\Pi_N\phi\).

All entries in (MR.8) are defined by the actual explicit-formula
Weil pairings.  They use only finitely many prime powers:
all arguments have support in
\([-R_N',R_N']\), where \(R_N'=\max(N,\sqrt2)\), so prime powers
larger than \(e^{2R_N'}\) make no contribution.  The remaining
archimedean and polar expressions are elementary computable
integrals of piecewise-polynomial functions and the explicit
archimedean kernel.  No zero ordinate or multiplicity is supplied
to these matrices.

\paragraph{A finite algorithm without exact real rounding.}
Set
\[
 \varepsilon_N=\frac{2^{-N}}{(d_N+1)^2}.
 \tag{MR.9}
\]
Compute a rational Hermitian matrix \(\widetilde S_N\), with
rational real and imaginary parts, and a rational real even
vector \(\widetilde a_N\), such that
\[
 \|\widetilde S_N-S_N\|_{\rm op}\le\varepsilon_N,\qquad
 \|\widetilde a_N-a_N\|_2\le\varepsilon_N,\qquad
 R_N\widetilde S_NR_N=-\widetilde S_N .
 \tag{MR.10}
\]
Choose these by one fixed precision and rational approximation
rule.  Symmetrization, followed by the odd parity projection,
preserves the required accuracy after starting at a smaller
fixed tolerance.

List the distinct positive eigenvalues of \(\widetilde S_N\) as
\(\theta_{N,1}>\cdots>\theta_{N,r_N}>0\), and put
\(t_{N,j}=1/\theta_{N,j}\).  Zero eigenvalues are discarded.
Use the full orthogonal eigenspace projection
\(E_{N,j}\), not a choice of eigenvector, to define
\[
 w_{N,j}=
 \frac{\|E_{N,j}\widetilde a_N\|_2^2}
      {\widehat\phi(t_{N,j})^2}.
 \tag{MR.11}
\]
Compute a nonnegative rational number \(u_{N,j}\) satisfying
\(|u_{N,j}-w_{N,j}|\le\varepsilon_N\).
Also compute a rational \(v_{N,j}\) in
\[
 [\,100\log(2+t_{N,j}),\,100\log(2+t_{N,j})+1\,].
\]
Define integers
\[
 k_{N,j}=\min\left\{
       \left\lfloor u_{N,j}+\tfrac12\right\rfloor,\,
       \left\lceil v_{N,j}\right\rceil+1
                 \right\},
 \qquad
 P_N(z)=\prod_{j=1}^{r_N}
             (1-\theta_{N,j}^2z^2)^{k_{N,j}} .
 \tag{MR.12}
\]
An exponent zero omits that factor.  This is a real polynomial
with real algebraic coefficients and only real zeros, normalized
by \(P_N(0)=1\).

\begin{theorem}[AC-S50: finite arithmetic reconstruction with multiplicities]
\label{AC-S50}
Assume \(A\).  The finite construction
(MR.4)--(MR.12) terminates for each \(N\), and
\[
                 P_N(z)\longrightarrow\frac{\Xi(z)}{\Xi(0)}
            \quad\hbox{locally uniformly on }\mathbb C .
 \tag{MR.13}
\]
For every fixed positive zero ordinate \(\gamma_j\), the finite
positive spectral point converging to it eventually receives
exactly its multiplicity \(m_{\gamma_j}\).
No simplicity assumption on the zeta zeros is made.
The assertion is about this fixed construction in the unshifted
Weil metric, not about the original lowest-eigenvector family.
\end{theorem}

\begin{proof}
\proofstep{0. The global sampling completion is onto.}
The spectral explicit formula makes the sampling map on
\(C_c^\infty(\mathbb R)\) an isometry for \(q\) into
\(\ell^2(\Gamma,m)\).  We prove that its image is dense.  Suppose
\(v=(v_\gamma)\in\ell^2(\Gamma,m)\) is orthogonal to this image.
The atomic expression
\[
 T=\sum_{\gamma\in\Gamma}m_\gamma\overline{v_\gamma}\,
                                      \delta_\gamma
\]
defines a tempered distribution.  Indeed Cauchy--Schwarz gives,
for every Schwartz function \(\varphi\),
\[
 |T(\varphi)|\le\|v\|_{\ell^2(\Gamma,m)}
       \left(\sum_\gamma m_\gamma|\varphi(\gamma)|^2\right)^{1/2},
\]
and the second factor is bounded by one fixed Schwartz seminorm,
using the multiplicity-counting estimate proved in Step 3 below.
Orthogonality says \(T(\widehat f)=0\) for every
\(f\in C_c^\infty\).  Compact smooth functions are dense in
Schwartz space, by multiplying a Schwartz function by expanding
smooth cutoffs.  The Fourier transform is a homeomorphism of
that space, so \(\mathcal F(C_c^\infty)\) is also dense in it.
Continuity therefore gives \(T=0\).  A compact test supported
near one discrete ordinate, equal to one there and vanishing
at the other ordinates, now gives
\(m_\gamma\overline{v_\gamma}=0\).  Thus \(v=0\).
The isometric completion map is consequently onto, establishing
(MR.2).  Multiplication by \(\gamma^{-1}\) is bounded, compact
and selfadjoint on this sequence space because zero is not an
ordinate and \(|\gamma|\to\infty\).  This proves (MR.3) and its
stated properties directly.

\proofstep{1. Density of the actual mean-zero spaces.}
At a fixed compact support the Weil form is continuous for the
\(H^1\) norm.  This follows directly from its explicit formula:
near zero the archimedean difference is bounded by
\[
 \|(T_h-1)f\|_2\|(T_h-1)g\|_2
       \le h^2\|f'\|_2\|g'\|_2,
\]
while the kernel is \(O(1/|h|)\); away from zero the bounded
support, finite prime sum and polar term give ordinary \(L^2\)
bounds.  Therefore rational dyadic piecewise-linear interpolation
is dense in compact smooth functions in the local form norm.
For a mean-zero smooth function, subtracting the integral error
times the fixed unit-integral hat \((1-|x|)_+\) makes every
approximant exactly mean zero, without changing convergence.
This proves density of \(\bigcup_NV_N\) in the global closure
of compact smooth mean-zero functions.

That latter closure is all of \(H\), by an explicit supported
correction.  Fix \(h\in C_c^\infty(\mathbb R)\) with
\(\int h=1\), and set \(h_R(x)=R^{-1}h(x/R)\), for integers
\(R\ge1\).  These functions are compactly supported, have
integral one, and integration by parts gives
\[
 |\widehat h_R(\gamma)|=|\widehat h(R\gamma)|
       \le \frac{\|h'\|_1}{R|\gamma|}.
\]
Since \(\sum_{\gamma\in\Gamma}m_\gamma/\gamma^2<\infty\),
\[
 q(h_R)\le \frac{\|h'\|_1^2}{R^2}
             \sum_{\gamma\in\Gamma}\frac{m_\gamma}{\gamma^2}
       \longrightarrow0.
 \tag{MR.13a}
\]
The summability follows from the order-one zero count, or directly
from the multiplicity-counting disk estimate in Step 3 below; it
is also the explicit input in Suzuki (2.2).  Given compact smooth
\(f\), the function \(f-(\int f)h_R\) is compact smooth and mean
zero, and its form distance from \(f\) tends to zero.  Since
compact smooth functions define the global completion, this
proves the claimed density without a prolate trial or a uniform
bound for integration.  Thus the projections \(\Pi_N\) onto
\(V_N\) tend strongly to \(I_H\).

\proofstep{2. Compact compression and the simple limiting coordinates.}
Compactness of \(C\) gives
\[
 \|\Pi_NC\Pi_N-C\|_{\rm op}\longrightarrow0.
 \tag{MR.14}
\]
For example approximate \(C\) in norm by a finite-rank operator;
strong convergence of \(\Pi_N\) is uniform on that finite-dimensional
range and on the adjoint range.  The rational approximation in
(MR.10), transported by the isometry following (MR.8) and extended
by zero on \(V_N^\perp\), differs from \(\Pi_NC\Pi_N\) in norm by at
most \(\varepsilon_N\).  It therefore also tends to \(C\) in norm.

For each fixed \(j\), \(1/\gamma_j\) is a simple isolated
eigenvalue of \(C\), even if the zeta zero has multiplicity.
A small circle about this eigenvalue excludes the rest of the
spectrum and zero.  The resolvent identity and a Neumann series
give convergence of the contour spectral projections.  Their
ranks are eventually one: two projections at distance less than
one have isomorphic ranges.  Hence the corresponding finite
positive eigenvalue is eventually simple and tends to
\(1/\gamma_j\).  Its eigenspace projection tends in norm to the
one-dimensional global projection.

If \(e_\gamma\) denotes the Weil-unit vector whose only nonzero
sampling coordinate is \(1/\sqrt{m_\gamma}\) at \(\gamma\), then
\[
       |q(\phi,e_\gamma)|^2
                   =m_\gamma\,|\widehat\phi(\gamma)|^2.
 \tag{MR.15}
\]
Since \(\Pi_N\phi\to\phi\) in \(H\) and the probe-coordinate error
tends to zero, (MR.15), the projection convergence and (MR.7)
show that \(w_{N,j}\to m_{\gamma_j}\), with the natural fixed-index
labeling for all sufficiently large \(N\).  This argument does
not evaluate a varying reproducing kernel at a moving spectral
point: it uses one fixed compact probe and its continuous,
strictly positive transform.

Finite eigenvalue degeneracy causes no ambiguity.  Formula
(MR.11) uses the whole spectral projection.  For each fixed
limiting coordinate the degeneracy eventually disappears by
the preceding rank argument.  Reflection pairs the positive
and negative spectral projections, and the evenness of the
probe makes their weights equal.

\proofstep{3. A fixed explicit cap.}
We verify a deliberately loose multiplicity majorant, so the
cap in (MR.12) uses no unknown zeta data.  Put
\(F(s)=(s-1)\zeta(s)\), which is entire.  For \(T\ge0\), use
the circle of radius \(3\) about \(2+iT\) and
\(M=\lceil T+8\rceil\).  On this circle
\(-1\le\operatorname{Re}s\le5\) and \(|s|\le M\).
The Euler--Maclaurin identity, DLMF 25.2.9 with \(n=1\), is
\[
 \zeta(s)=\sum_{n=1}^M n^{-s}
       +\frac{M^{1-s}}{s-1}
       -\frac12M^{-s}+\frac{s}{12}M^{-s-1}
       -\frac{s(s+1)(s+2)}6
           \int_M^\infty\frac{\widetilde B_3(x)}{x^{s+3}}\,dx .
 \tag{MR.16}
\]
It holds for \(\operatorname{Re}s>-2\).
Since \(B_3(x)=x(x-\tfrac12)(x-1)\) on \([0,1]\),
\(\|\widetilde B_3\|_\infty\le1\).  The last term has modulus at
most \((M+1)(M+2)/6\); the preceding terms, after multiplication
by \(s-1\), give the elementary bound
\[
  |F(s)|\le10(M+3)^3
       \le17280(1+T)^3.
 \tag{MR.17}
\]
This loose bound follows term by term from
\(\sum_{n\le M}n\le M(M+1)/2\),
\(|s-1|\le M+1\), and
\(\int_M^\infty x^{-\operatorname{Re}s-3}dx
 \le M^{-1}\).
The apparent pole at \(s=1\) is removable in \(F\).

At the center the Euler product gives
\[
 |F(2+iT)|\ge\frac{\sqrt{1+T^2}}{\zeta(2)}
             \ge\frac{1+T}{3}.
\]
Jensen's formula therefore bounds the sum of the zero contributions
in that disk by
\[
              \log51840+2\log(1+T).
 \tag{MR.18}
\]
A critical-line zero at ordinate \(T\) has distance \(3/2\)
from the center and contributes its multiplicity times \(\log2\).
Consequently
\[
 m_T\le\frac{\log51840+2\log(1+T)}{\log2}
                         <100\log(2+T).
 \tag{MR.19}
\]
The same Jensen bound, using radius \(2\) inside the radius-\(3\)
disk, counts every critical-line zero with ordinate in
\([T-1,T+1]\); their distances are at most \(\sqrt{13}/2<2\).
It gives \(O(\log(2+T))\) zeros in that interval, with
multiplicity.  Thus
\[
 \sum_{j\ge1}\frac{\log(2+\gamma_j)}{\gamma_j^2}<\infty.
 \tag{MR.20}
\]
The finitely many small ordinates cause no problem because zero
is not an ordinate.  One may also obtain (MR.20) from the usual
zero-counting theorem; the elementary disk argument fixes all
that is needed here.

For each fixed \(j\), (MR.19) and convergence of \(t_{N,j}\)
give a strict margin below the cap eventually.  Also
\(u_{N,j}\to m_{\gamma_j}\).  Eventually it is within \(1/4\)
of that integer, so rounding is unambiguous and
\[
                     k_{N,j}=m_{\gamma_j}.
 \tag{MR.21}
\]
No test for equality of an exact real weight to a half-integer
has been required.

\proofstep{4. Uniform tails, including finite numerical perturbations.}
Order the distinct positive finite eigenvalues as in (MR.11).
The \(j\)-th distinct one occurs no earlier than the \(j\)-th
positive eigenvalue counted with multiplicity.  The min--max
principle, followed by the operator-norm error (MR.10), yields
\[
 \theta_{N,j}\le\frac1{\gamma_j}+\varepsilon_N,\qquad
 t_{N,j}\ge\left(\gamma_j^{-1}+\varepsilon_N\right)^{-1}.
 \tag{MR.22}
\]
This bound also covers small positive eigenvalues created by
the rational perturbation of a zero eigenvalue.  It asserts no
monotonicity for the perturbed matrices.

The cap gives
\(0\le k_{N,j}\le100\log(2+t_{N,j})+3
 \le105\log(2+t_{N,j})\).
The function \(g(t)=\log(2+t)/t^2\) decreases for \(t>0\).
Split (MR.22) into \(\gamma_j\le\varepsilon_N^{-1}\) and
\(\gamma_j>\varepsilon_N^{-1}\).  In the first case
\(t_{N,j}\ge\gamma_j/2\); in the second
\(t_{N,j}\ge(2\varepsilon_N)^{-1}\).  Hence, for every \(J\),
\[
 \sum_{\substack{j>J\\j\le r_N}}
       \frac{\log(2+t_{N,j})}{t_{N,j}^2}
 \le4\sum_{j>J}\frac{\log(2+\gamma_j)}{\gamma_j^2}
   +4d_N\varepsilon_N^2
                       \log(2+(2\varepsilon_N)^{-1}).
 \tag{MR.23}
\]
The last term tends to zero by the explicit (MR.4), (MR.9).
For a compact disk \(|z|\le R\), choose \(J\) so large that
\(\gamma_{J+1}/2>2R\), and then take \(N\) sufficiently large.
All factors in the displayed tail have \(|z^2/t_{N,j}^2|<1/4\).
The analytic logarithm based at \(1\) obeys
\[
 \left|\sum_{j>J}k_{N,j}
        \log(1-z^2/t_{N,j}^2)\right|
       \le210R^2\sum_{j>J}g(t_{N,j}).
 \tag{MR.24}
\]
Equations (MR.20), (MR.23) make this uniformly small, first for
large \(J\) and then large \(N\).  The fixed first \(J\) factors
converge by (MR.21) and spectral convergence.  Therefore
\[
 P_N(z)\longrightarrow
       \prod_{j\ge1}(1-z^2/\gamma_j^2)^{m_{\gamma_j}}
 \quad\hbox{locally uniformly}.
\]
The usual order-one Hadamard product for the entire even
\(\Xi\), with its multiplicities, identifies this product with
\(\Xi(z)/\Xi(0)\).  Pairing the factors at \(\gamma\) and
\(-\gamma\) cancels their genus-one exponential factors; evenness
removes the remaining linear exponential, and normalization at
zero removes the constant.  This proves (MR.13).

\proofstep{5. What arithmetic computability means here.}
All basis functions and their primitives are explicitly specified.
The explicit formula computes every exact pairing in (MR.8) to
certified arbitrary precision from a finite list of prime powers,
the polar term and archimedean integrals.  The near-zero
archimedean cancellation estimated in Step 1 bounds the omitted
small integration interval, and the exponential tail of that
kernel bounds the omitted infinite interval.  On the remaining
compact subintervals the piecewise-polynomial expressions give
ordinary effective quadrature bounds.

Under \(A\), positivity of \(G_N\) lets an increasing-precision
search certify a positive lower eigenvalue bound and hence compute
its positive inverse square root and (MR.8) to the accuracy
(MR.10).  No uniform lower bound on this Gram matrix is assumed;
the finite search may be expensive but terminates.
Rational Hermitian \(\widetilde S_N\) has algebraic eigenvalues.
Their signs, exact multiplicities, distinct eigenspaces and
orthogonal projections can be determined by exact algebraic
root isolation and linear algebra.  Thus exact finite degeneracy
and zero eigenvalues require no oracle for equality of arbitrary
computable real numbers.

The denominator in (MR.11) is a known strictly positive computable
real.  An increasing-precision search bounds it away from zero
for each fixed algebraic argument and permits a rational weight
approximation of the prescribed accuracy.  Rounding is then done
on that rational number; either convention at a rational tie is
fixed once and for all.  The logarithmic cap is also chosen from
a rational enclosure, so exact logarithmic integer-boundary tests
are unnecessary.  The construction therefore uses no zeta-zero
location, zero derivative or multiplicity as finite input, and
does not obtain its coefficients from an unevaluated limiting
operator.  It claims neither a practical complexity bound nor
an effective convergence modulus in (MR.13).
\end{proof}

\paragraph{Source scope and boundary.}
The complete printed Suzuki paper, including its inverse-compression,
compact-convergence and product arguments, was read in the preceding
exact-scope audit.  The present proof supplies sampling surjectivity
in Step 0 rather than importing its cited global realization theorem.
The ordinary explicit formula and local strict positivity already
recorded in AC-S03 remain background inputs.  The new finite-space
density, compact primitive matrix, fixed-probe multiplicity extraction,
rational perturbation, integer-rounding and capped-tail arguments
are proved above.
The applied DLMF Euler--Maclaurin formula, its stated half-plane,
Euler product and Bernoulli definitions were checked directly;
no full proof of every cited handbook or de Branges input is claimed.

The Gram pairings are genuine finite-prime arithmetic data, and the
finite polynomials have real zeros.  Their construction is fixed after
the explicitly stated dyadic-basis and rational-approximation rules.
It is not the S25 determinant of the original ground line, not a new
motivic or etale cohomology object, and not a solution of the remaining
positive spectral-edge comparison.  The latter still requires its
own quantitative control.
